\%============================================================
\chapter{Bài Học: Kiến Tạo Cuộc Đời --- Bạn Là Kiến Trúc Sư Của Chính Mình (Building Your Life --- You Are Your Own Architect)}
\begin{center}
  \textit{\color{truyengold}``The future belongs to those who believe in the beauty of their dreams.''}\\[4pt]
  \textit{(Tương lai thuộc về những ai tin vào vẻ đẹp của ước mơ của họ.)}\\[4pt]
  \small\color{truyenblue}--- Cổ Kim Đông Tây ---
\end{center}
\ngancach

\section{Steve Jobs Và Nghệ Thuật Kết Nối Những Dấu Chấm}
\begin{truyen}{Steve Jobs Và Nghệ Thuật Kết Nối Những Dấu Chấm}{Steve Jobs --- Hoa Kỳ, 2005}
\chuhoa{S}{teve} Jobs trong bài phát biểu tốt nghiệp nổi tiếng tại Stanford năm 2005 kể câu chuyện về những dấu chấm.\\[4pt]
\textit{(Steve Jobs in his famous 2005 Stanford commencement address told the story of connecting the dots.)}

Sau khi bỏ học đại học, ông vẫn ở lại nghe lén các lớp học --- và học thêm khoá thư pháp vì cảm thấy thú vị.\\[4pt]
\textit{(After dropping out of college, he still stayed and sneaked into classes --- and studied calligraphy because he found it fascinating.)}

Mười năm sau, khi thiết kế Mac, ông đưa thư pháp vào máy tính --- tạo ra những font chữ đẹp đầu tiên trong lịch sử.\\[4pt]
\textit{(Ten years later, designing the Mac, he put calligraphy into the computer --- creating the first beautiful typefaces in history.)}

Ông nói: ``Bạn không thể kết nối những dấu chấm nhìn về phía trước; bạn chỉ có thể kết nối nhìn ngược lại.''\\[4pt]
\textit{(He said: `You cannot connect the dots looking forward; you can only connect them looking backwards.')}

``Vì vậy bạn phải tin rằng những dấu chấm sẽ được kết nối theo một cách nào đó trong tương lai.''\\[4pt]
\textit{(`So you must trust that the dots will somehow connect in your future.')}

\end{truyen}
\begin{baihoc}
  Mỗi điều bạn học --- dù có vẻ vô dụng tức thì --- đều có thể trở thành viên gạch xây nên kỳ tích của bạn về sau.\\[4pt]
  \textit{(Every thing you learn --- however apparently useless immediately --- can become a brick building your future masterpiece.)}
\end{baihoc}

\section{Walt Disney --- Bảy Lần Phá Sản, Một Lần Đế Chế}
\begin{truyen}{Walt Disney --- Bảy Lần Phá Sản, Một Lần Đế Chế}{Walt Disney --- Hoa Kỳ, 1966}
\chuhoa{W}{alt} Disney bị đuổi khỏi tờ báo vì ``thiếu trí tưởng tượng và không có ý tưởng tốt nào.''\\[4pt]
\textit{(Walt Disney was fired from a newspaper for `lacking imagination and having no good ideas.')}

Ông thành lập công ty hoạt hình đầu tiên và phá sản; thành lập lại và phá sản lần hai.\\[4pt]
\textit{(He founded his first animation company and went bankrupt; founded another and went bankrupt again.)}

Mickey Mouse bị 300 nhà đầu tư từ chối trước khi một người chịu nghe --- và khi phim Bạch Tuyết ra mắt, mọi người kinh ngạc.\\[4pt]
\textit{(Mickey Mouse was rejected by 300 investors before one listened --- and when Snow White premiered, audiences were astonished.)}

Disney nói: ``Tất cả những ước mơ của chúng ta có thể trở thành sự thật nếu chúng ta có can đảm theo đuổi chúng.''\\[4pt]
\textit{(Disney said: `All our dreams can come true if we have the courage to pursue them.')}

Đế chế Disney ngày nay là một trong những thương hiệu giá trị nhất thế giới --- sinh ra từ những lần phá sản và những giấc mơ không tắt.\\[4pt]
\textit{(The Disney empire today is one of the world's most valuable brands --- born from bankruptcies and undying dreams.)}

\end{truyen}
\begin{baihoc}
  Giấc mơ không chết vì thất bại --- nó chỉ chết khi bạn quyết định ngừng theo đuổi nó.\\[4pt]
  \textit{(A dream does not die from failure --- it only dies when you decide to stop pursuing it.)}
\end{baihoc}

\section{Oprah Winfrey --- Từ Nghèo Khổ Đến Ảnh Hưởng}
\begin{truyen}{Oprah Winfrey --- Từ Nghèo Khổ Đến Ảnh Hưởng}{Oprah Winfrey --- Hoa Kỳ, 1954}
\chuhoa{O}{prah} Winfrey sinh ra trong nghèo khổ ở Mississippi; bị lạm dụng từ nhỏ và mang thai năm 14 tuổi.\\[4pt]
\textit{(Oprah Winfrey was born into poverty in Mississippi; she was abused as a child and became pregnant at 14.)}

Bị đuổi khỏi chương trình tin tức vì ``quá cảm xúc''; đã tưởng sự nghiệp truyền hình kết thúc.\\[4pt]
\textit{(She was fired from a news programme for being `too emotional'; her television career seemed finished.)}

Nhưng bà tiếp tục --- và show Oprah Winfrey trở thành chương trình talkshow thành công nhất lịch sử truyền hình Mỹ.\\[4pt]
\textit{(But she continued --- and The Oprah Winfrey Show became the most successful talk show in American television history.)}

Bà nói: ``Không có người thành công nào mà không từng thất bại nhiều lần; điều khác biệt là họ không đứng yên sau khi ngã.''\\[4pt]
\textit{(She said: `No truly successful person has not failed many times; what is different is they don't stay down after falling.')}

Bà dùng tầm ảnh hưởng của mình để trao học bổng, thành lập trường học ở Nam Phi và cổ suý đọc sách trên toàn cầu.\\[4pt]
\textit{(She used her influence to fund scholarships, establish a school in South Africa, and champion reading globally.)}

\end{truyen}
\begin{baihoc}
  Câu chuyện bạn kể về chính mình quyết định tất cả --- ai nói mình không thể thì thường đúng về bản thân họ.\\[4pt]
  \textit{(The story you tell about yourself determines everything --- those who say they cannot are usually right about themselves.)}
\end{baihoc}

\section{Thomas Edison Và Một Nghìn Bước Tới Thành Công}
\begin{truyen}{Thomas Edison Và Một Nghìn Bước Tới Thành Công}{Thomas Edison --- Hoa Kỳ, 1879}
\chuhoa{T}{homas} Edison thử nghiệm gần một nghìn loại vật liệu trước khi tìm ra sợi than cho bóng đèn điện.\\[4pt]
\textit{(Thomas Edison tested nearly one thousand materials before finding the carbon filament for the electric light bulb.)}

Khi phóng viên hỏi cảm giác thất bại cả nghìn lần, ông trả lời: ``Tôi chưa thất bại --- tôi đã học 999 cách không hoạt động.''\\[4pt]
\textit{(When a reporter asked how it felt to fail a thousand times, he replied: `I have not failed --- I have learned 999 ways that do not work.')}

Edison nói: ``Thiên tài là một phần trăm cảm hứng và chín mươi chín phần trăm mồ hôi và nỗ lực.''\\[4pt]
\textit{(Edison said: `Genius is one percent inspiration and ninety-nine percent perspiration.')}

Phòng thí nghiệm của ông ở Menlo Park đăng ký hơn một nghìn bằng sáng chế --- trung bình mười bốn ngày một bằng.\\[4pt]
\textit{(His laboratory at Menlo Park registered over one thousand patents --- an average of one every fourteen days.)}

Chìa khoá không phải là thiên tài hay may mắn mà là hệ thống làm việc có kỷ luật và không chịu bỏ cuộc.\\[4pt]
\textit{(The key was not genius or luck but a disciplined work system and refusing to give up.)}

\end{truyen}
\begin{baihoc}
  Thành công không đến một lần --- nó là kết quả tích luỹ của hàng trăm lần thử, học và thử lại.\\[4pt]
  \textit{(Success does not come once --- it is the cumulative result of hundreds of attempts, learnings, and attempts again.)}
\end{baihoc}

\section{J.K. Rowling --- Hành Trình Từ Vô Danh Đến Huyền Thoại}
\begin{truyen}{J.K. Rowling --- Hành Trình Từ Vô Danh Đến Huyền Thoại}{J.K. Rowling --- Anh, 1997}
\chuhoa{J}{oanne} Rowling viết Harry Potter trên những tờ giấy ăn trong quán cà phê ở Edinburgh --- người mẹ đơn thân đang sống nhờ trợ cấp.\\[4pt]
\textit{(Joanne Rowling wrote Harry Potter on napkins in an Edinburgh café --- a single mother living on welfare.)}

Mười hai nhà xuất bản từ chối bản thảo; nhà xuất bản thứ mười ba nhận với điều kiện bà không hy vọng kiếm tiền từ sách thiếu nhi.\\[4pt]
\textit{(Twelve publishers rejected the manuscript; the thirteenth accepted on condition she not expect to earn money from children's books.)}

Harry Potter và Hòn Đá Phù Thuỷ ra mắt năm 1997 và trở thành cơn sốt toàn cầu.\\[4pt]
\textit{(Harry Potter and the Philosopher's Stone was published in 1997 and became a worldwide phenomenon.)}

Rowling nói: ``Bạn không thể nhìn thấy con đường phía trước từ nơi bạn đang đứng; bạn chỉ thấy từng bước tiếp theo.''\\[4pt]
\textit{(Rowling said: `You cannot see the path ahead from where you stand; you can only see the very next step.')}

Cuộc sống của bà là bằng chứng rằng thất bại không phải điểm cuối --- nó là phần không thể thiếu của hành trình.\\[4pt]
\textit{(Her life is proof that failure is not the end --- it is an essential part of the journey.)}

\end{truyen}
\begin{baihoc}
  Bước tiếp theo --- không phải cánh chung --- là tất cả những gì bạn cần nhìn thấy để tiếp tục đi.\\[4pt]
  \textit{(The very next step --- not the final destination --- is all you need to see in order to keep going.)}
\end{baihoc}

\section{Elon Musk Và Những Canh Bạc Không Ai Dám Làm}
\begin{truyen}{Elon Musk Và Những Canh Bạc Không Ai Dám Làm}{Elon Musk --- Hoa Kỳ, 2002}
\chuhoa{E}{lon} Musk sau khi bán PayPal có 180 triệu đô la và đặt câu hỏi: điều gì quan trọng nhất cho tương lai của nhân loại?\\[4pt]
\textit{(Elon Musk after selling PayPal had 180 million dollars and asked: what matters most for humanity's future?)}

Câu trả lời của ông: năng lượng bền vững và thuộc địa hoá đa hành tinh để bảo đảm sự tồn vong của loài người.\\[4pt]
\textit{(His answer: sustainable energy and multi-planet colonisation to ensure humanity's survival.)}

Ông đặt cược phần lớn tài sản vào Tesla và SpaceX --- cả hai đều suýt phá sản nhiều lần.\\[4pt]
\textit{(He bet most of his wealth on Tesla and SpaceX --- both came near bankruptcy multiple times.)}

Năm 2008, trong cuộc khủng hoảng tài chính, ông hầu như không còn tiền và mượn bạn bè để trả lương nhân viên.\\[4pt]
\textit{(In 2008, during the financial crisis, he was almost out of money and borrowed from friends to pay employees.)}

SpaceX trở thành công ty tư nhân đầu tiên đưa tàu vũ trụ lên quỹ đạo; Tesla trở thành hãng xe điện giá trị nhất thế giới.\\[4pt]
\textit{(SpaceX became the first private company to send spacecraft to orbit; Tesla became the world's most valuable electric car company.)}

\end{truyen}
\begin{baihoc}
  Những ước mơ lớn đòi hỏi những canh bạc lớn --- và chỉ những người dám đặt cược bản thân mới có thể thay đổi thế giới.\\[4pt]
  \textit{(Big dreams require big bets --- and only those willing to wager themselves can change the world.)}
\end{baihoc}

\section{Warren Buffett Và Triết Học Đầu Tư Dài Hạn}
\begin{truyen}{Warren Buffett Và Triết Học Đầu Tư Dài Hạn}{Warren Buffett --- Hoa Kỳ, 1956}
\chuhoa{W}{arren} Buffett mua cổ phiếu đầu tiên năm 11 tuổi và học bài học đắt giá đầu tiên khi nó giảm giá.\\[4pt]
\textit{(Warren Buffett bought his first stock at age 11 and learned his first costly lesson when it fell in value.)}

Ông đọc hơn 500 trang sách và báo cáo kinh doanh mỗi ngày --- hành động này ông duy trì suốt hơn 70 năm đầu tư.\\[4pt]
\textit{(He reads more than 500 pages of books and business reports every day --- a practice he has maintained throughout over 70 years of investing.)}

Triết học đầu tư của ông: ``Hãy mua cổ phiếu khi người khác sợ hãi và bán khi người khác tham lam.''\\[4pt]
\textit{(His investment philosophy: `Buy when others are fearful and sell when others are greedy.')}

Ông nói: ``Ai ngồi mát mà không gieo hạt sẽ không có gì mà gặt khi mùa thu đến.''\\[4pt]
\textit{(He says: `One who sits in shade without planting will have nothing to harvest when autumn comes.')}

Buffett cam kết tặng 99 phần trăm tài sản của mình cho từ thiện --- đó là triết học sống, không phải chiến lược PR.\\[4pt]
\textit{(Buffett has pledged to give 99 percent of his wealth to charity --- that is a life philosophy, not a PR strategy.)}

\end{truyen}
\begin{baihoc}
  Kiến tạo cuộc đời là hành động lâu dài --- mỗi ngày gieo một hạt giống nhỏ vào đúng đất là cách thay đổi bền vững nhất.\\[4pt]
  \textit{(Building a life is a long-term endeavour --- planting one small seed daily in the right soil is the most sustainable way to transform.)}
\end{baihoc}

\section{Hồ Chí Minh Và Ý Chí Kiến Tạo Đất Nước}
\begin{truyen}{Hồ Chí Minh Và Ý Chí Kiến Tạo Đất Nước}{Hồ Chí Minh --- Việt Nam, 1945}
\chuhoa{N}{ăm} 1945, Hồ Chí Minh đọc Tuyên Ngôn Độc Lập của Việt Nam --- bắt đầu bằng trích dẫn từ Tuyên Ngôn Độc Lập Mỹ.\\[4pt]
\textit{(In 1945, Ho Chi Minh read Vietnam's Declaration of Independence --- beginning with a quote from the American Declaration of Independence.)}

Đó là tuyên bố về tầm nhìn: Việt Nam sẽ xây dựng một đất nước bình đẳng và tự do, như những gì nhân loại mơ ước.\\[4pt]
\textit{(It was a vision statement: Vietnam would build a nation of equality and freedom, as all humanity dreams.)}

Ông đã dành 30 năm bôn ba khắp thế giới --- Paris, Moscow, Trung Quốc, Thái Lan --- học hỏi và chuẩn bị.\\[4pt]
\textit{(He had spent 30 years travelling the world --- Paris, Moscow, China, Thailand --- learning and preparing.)}

Khi về nước lãnh đạo, ông không chỉ là nhà cách mạng mà là người thầy --- dạy dân đọc, dạy dân vệ sinh, dạy đạo đức.\\[4pt]
\textit{(When he returned to lead the country, he was not just a revolutionary but a teacher --- teaching literacy, hygiene, and ethics.)}

Di sản của ông là một nước Việt Nam có độc lập --- và hành trình xây dựng đó vẫn tiếp tục sau ông.\\[4pt]
\textit{(His legacy is an independent Vietnam --- and the building journey he began continues after him.)}

\end{truyen}
\begin{baihoc}
  Kiến tạo đất nước cũng như kiến tạo cuộc đời --- cần tầm nhìn rõ ràng, sự chuẩn bị kỹ lưỡng và ý chí không khuất phục.\\[4pt]
  \textit{(Building a nation is like building a life --- it needs clear vision, thorough preparation, and an indomitable will.)}
\end{baihoc}

\section{Marie Curie Và Sự Kiên Trì Trong Khoa Học}
\begin{truyen}{Marie Curie Và Sự Kiên Trì Trong Khoa Học}{Marie Curie --- Pháp, 1898}
\chuhoa{M}{arie} Curie là người phụ nữ đầu tiên nhận giải Nobel và người đầu tiên trong lịch sử nhận giải Nobel hai lần ở hai lĩnh vực khác nhau.\\[4pt]
\textit{(Marie Curie was the first woman to receive a Nobel Prize and the first person ever to receive two Nobels in two different fields.)}

Sinh ra ở Ba Lan khi phụ nữ không được vào đại học, bà đến Paris học lậu và sau đó học chính thức.\\[4pt]
\textit{(Born in Poland where women could not attend university, she went to Paris to study and eventually gained formal admission.)}

Bà và Pierre làm việc trong điều kiện cực kỳ thiếu thốn --- phòng thí nghiệm bị dột, mùa đông lạnh giá không có sưởi.\\[4pt]
\textit{(She and Pierre worked in extremely poor conditions --- a leaking laboratory, freezing winters without heat.)}

Bà xử lý hàng tấn quặng uranium bằng tay --- không biết rằng phóng xạ đang từ từ giết chết bà.\\[4pt]
\textit{(She processed tons of uranium ore by hand --- not knowing that radiation was slowly killing her.)}

Bà nói: ``Trong cuộc sống không có gì đáng sợ, chỉ có những điều cần được hiểu. Giờ là lúc hiểu thêm để sợ ít hơn.''\\[4pt]
\textit{(She said: `Nothing in life is to be feared, only to be understood. Now is the time to understand more, so that we may fear less.')}

\end{truyen}
\begin{baihoc}
  Hành trình khoa học và hành trình sống giống nhau: bạn không thấy hết trước hành trình; nhưng mỗi bước sáng tỏ thêm một chút.\\[4pt]
  \textit{(The scientific journey and life journey are alike: you cannot see the whole path ahead; but each step illuminates a little more.)}
\end{baihoc}

\section{Nguyễn Trường Tộ Và Giấc Mơ Canh Tân Dang Dở}
\begin{truyen}{Nguyễn Trường Tộ Và Giấc Mơ Canh Tân Dang Dở}{Nguyễn Trường Tộ --- Việt Nam, thế kỷ 19}
\chuhoa{N}{guyễn} Trường Tộ là nhà trí thức Công giáo người Việt sống cuối thế kỷ 19, đi châu Âu và thấy sự tiến bộ của phương Tây.\\[4pt]
\textit{(Nguyen Truong To was a Vietnamese Catholic intellectual of the late 19th century who went to Europe and saw Western progress.)}

Ông viết hàng chục bản điều trần gửi vua Tự Đức và triều đình, đề nghị canh tân đất nước toàn diện.\\[4pt]
\textit{(He wrote dozens of policy memoranda to Emperor Tu Duc and the court, proposing comprehensive modernisation of the country.)}

Ông đề xuất cải cách giáo dục, học khoa học kỹ thuật phương Tây, mở cửa thương mại --- trước Nhật Bản cải cách Minh Trị.\\[4pt]
\textit{(He proposed education reform, learning Western science and technology, opening trade --- before Japan's Meiji reforms.)}

Triều đình không nghe; Pháp xâm lược và chiếm Việt Nam --- sau đó người ta mới tiếc những điều ông đã cảnh báo.\\[4pt]
\textit{(The court did not listen; France invaded and occupied Vietnam --- only then did people regret what he had warned.)}

Giấc mơ của ông dang dở --- nhưng tầm nhìn của ông vẫn là tấm gương cho những ai muốn kiến tạo tương lai của dân tộc.\\[4pt]
\textit{(His dream remained unfulfilled --- but his vision remains a mirror for all who wish to build their nation's future.)}

\end{truyen}
\begin{baihoc}
  Kiến tạo tương lai đòi hỏi can đảm nói sự thật cho người không muốn nghe --- dù kết quả không được thấy trong đời mình.\\[4pt]
  \textit{(Building the future requires the courage to tell truth to those who do not wish to hear --- even if you never see the result in your lifetime.)}
\end{baihoc}
